\chapter{Lecture Sources and Historical Notes}\label{ch:source-notes}
The course materials supplied with this project are the primary reference for its scope and terminology.
Slide numbers refer to the printed identifiers inside the handouts, which may differ from PDF page numbers because a PDF page can contain two slides.

\section*{How Attribution Works}
Lecture examples are paraphrased, not reconstructed quotations.
They preserve the teaching point while making modeling assumptions explicit.
Examples added to connect ideas, rather than recovered from a lecture, are presented as the book's own illustrations.

\section*{Examples Used in the Narrative}
\begin{description}
  \item[Independent additions, Chapter 1.] \texttt{n1\_intro.pdf}, INTRO-6, contrasts three dependent additions with three independent additions on three adders. The book states the equal-latency, available-resource assumptions in its diagram.
  \item[Scheduling the next crew, Chapter 1.] January 13, 2026 lecture, Notability note \emph{1. Syllabus}, \emph{Recording 2}, approximately 14:41--15:25. The construction example separates estimating a completion time from acting on that estimate.
  \item[The transistor-budget question, Chapter 1.] January 15, 2026 lecture, Notability note \emph{2. ECE 6100}, \emph{Recording 3}, approximately 06:47--08:05, and slides INTRO-18--19. The book retains the design question without treating the power wall as the end of single-core improvement.
  \item[The Cray vector example, Chapter 1.] \texttt{n1\_intro.pdf}, INTRO-13--14, contrasts streamed vector execution with packed-word arithmetic. The book separates the lecture's simplified pipeline from historical hardware timings.
  \item[Buying a jet-simulation machine, Chapter 2.] \texttt{n2\_perf.pdf}, PERF-18--20, uses an extravagant many-processor machine to expose Amdahl's limit, then changes the question to a larger simulation for Gustafson's law. The book uses 100 hours as a convenient baseline rather than the handout's rounded week-to-seconds conversion.
  \item[Passenger time and airline throughput, Chapter 2.] January 20, 2026 lecture, Notability note \emph{3. ECE 6100}, \emph{Recording 1}, approximately 02:00--05:15. The book keeps the qualitative Concorde-versus-larger-aircraft comparison and states its simplified capacity assumptions.
  \item[Buying the slower machine, Chapter 2.] \texttt{n2\_perf.pdf}, PERF-9--11, compares rates of 4 and 4 with rates of 2 and 7. The workload must contain equal work of each type for the harmonic-mean comparison.
  \item[The supercomputer joke, Chapter 2.] January 15, 2026 lecture, Notability note \emph{2. ECE 6100}, \emph{Recording 2}, approximately 06:03--06:32.
  \item[Choosing a victim, Chapter 3.] \texttt{n3\_cache\_all.pdf}, CACHE-20, asks which of four blocks should be replaced after a short access history. The book follows that example to distinguish knowing the past from knowing the next use.
  \item[Why renaming mattered, Chapter 6.] \texttt{n5\_ilp\_part1.pdf}, ILP-14, connects the System/360 Model 91's long floating-point operations and small architectural register set to dynamic scheduling and renaming.
\end{description}

\section*{Primary Historical References}
\begin{description}
  \item[Multicore and energy efficiency.] Intel's \href{https://www.intel.com/pressroom/archive/releases/2006/20060307corp.htm}{March 7, 2006 Core microarchitecture introduction} documents the contemporary emphasis on multicore performance and energy efficiency. It provides historical context for the transistor-budget discussion in Chapter 1.
  \item[Vector registers and chaining.] Cray Research, \emph{CRAY-1 Computer Systems Hardware Reference Manual}, HR-0004F, May 1982, sections on vector registers and chaining. The manual describes eight vector registers with up to 64 elements each and explains the conditions for forwarding an element stream between functional units. Available from the Computer History Museum's \href{https://bitsavers.computerhistory.org/pdf/cray/CRAY-1/HR-0004F_CRAY-1_Computer_Systems_Hardware_Reference_Manual_May82.pdf}{Bitsavers archive}.
  \item[Dynamic scheduling.] R. M. Tomasulo, \emph{An Efficient Algorithm for Exploiting Multiple Arithmetic Units}, IBM Journal of Research and Development, 11(1), 1967, pp.\ 25--33. This is the original description of the Model 91 floating-point scheduling method. \href{https://doi.org/10.1147/rd.111.0025}{doi:10.1147/rd.111.0025}.
  \item[Two-level predictor organizations.] Tse-Yu Yeh and Yale N. Patt, \emph{Alternative Implementations of Two-Level Adaptive Branch Prediction}, ISCA, 1992. The \href{https://people.eecs.berkeley.edu/~kubitron/courses/cs252-S09/handouts/papers/p451-yeh.pdf}{paper} distinguishes global and per-address histories and shared versus separate pattern tables. Chapter 4 identifies its local-history example within that larger family.
  \item[Perceptron predictors.] Daniel A. Jiménez and Calvin Lin, \emph{Dynamic Branch Prediction with Perceptrons}, HPCA, 2001. The \href{https://www.cs.utexas.edu/~lin/papers/hpca01.pdf}{authors' paper}, Sections 3 and 5, supplies the training rule, linear-separability limitation, and empirically tuned history-length threshold used in Chapter~4.
\end{description}

\section*{Reading the Remaining Course Handouts}
The chapter sequence follows \path{n1_intro.pdf}, \path{n2_perf.pdf}, \path{n3_cache_all.pdf}, \path{n4_pipe_part1.pdf}, \path{n5_ilp_part1.pdf}, the VLIW/software-pipelining handouts, \path{n6_multiproc v1.1.pdf}, \path{n7_synch.pdf}, and \path{n8_interconn.pdf}.
These filenames identify the local course sources; the handouts are not redistributed with the book.
Course-specific categories, including the broad use of speculation for locality-based caching, are identified in the text when a narrower definition is also common.

\section*{Specifications Behind Implementation Caveats}
\begin{description}
  \item[Atomic operations and reservations.] The \href{https://docs.riscv.org/reference/isa/unpriv/a-st-ext.html}{RISC-V atomic-extension specification} provides a concrete example of selectable acquire/release ordering and permitted store-conditional failure. These details motivate the distinction between atomicity and ordering in Chapter~\ref{ch:sync-consistency}; other ISAs have their own rules.
  \item[Dynamic and guided loop scheduling.] The \href{https://www.openmp.org/spec-html/5.2/openmpse66.html}{OpenMP worksharing-loop specification} distinguishes fixed-size dynamic chunks from guided chunks that decrease as work is consumed. Chapter~\ref{ch:sync-consistency} uses this distinction when describing a fetch-and-add work counter.
  \item[Persistence and cached writes.] Intel's \href{https://www.intel.com/content/www/us/en/developer/articles/technical/enabling-persistent-memory-in-the-storage-performance-development-kit-spdk.html}{persistent-memory programming discussion} explains why writes can remain in volatile caches or buffers and need appropriate propagation to reach a persistence domain. Chapter 3 separates this requirement from write-through policy and cache coherence.
  \item[Translation faults.] The \href{https://docs.riscv.org/reference/isa/priv/supervisor.html}{RISC-V supervisor-level specification}, under virtual-address translation, illustrates that page-table validity and permission checks can raise faults. Chapter 3 distinguishes a cached-translation miss from a page fault and from storage I/O.
\end{description}
